\chapter{Phương pháp đề xuất}
\label{chap:method}

% Cho phép co giãn khẩn cấp để ngắt dòng trước các định danh monospace dài; giới hạn trong chương.
\begingroup
\emergencystretch=3em

Phương pháp đề xuất phát hiện lỗ hổng phân quyền có trạng thái bằng cách đối chiếu kết quả của cùng một thao tác trước và sau khi quyền của chủ thể bị thu hồi. Chương này trình bày phương pháp từ quan hệ biến hình MR-A1 đến kịch bản kiểm thử chạy được trong EvoMaster.

\section{Tổng quan Oracle 970}
\label{sec:method_overview}

Trong kiểm thử phân quyền REST API, bài toán thiếu oracle (Mục~\ref{sec:bg_metamorphic}) xuất hiện vì mã trạng thái HTTP không cho biết thao tác có được phép theo chính sách hay không. Máy chủ có thể trả \texttt{200 OK} cho một thao tác mà chính sách phân quyền cấm.

Oracle 970 chạy trong pha kiểm thử bảo mật của EvoMaster 6.1.1 và dùng cùng kiến trúc với các oracle an toàn có sẵn của công cụ (Mục~\ref{sec:bg_related_work}). Pha này lấy đầu vào từ kho lưu trữ ca kiểm thử (Mục~\ref{sec:bg_sbst}). Số 970 là mã lỗi phần mở rộng gán cho vi phạm MR-A1: EvoMaster 6.1.1 dành dải mã 9xx cho các oracle thử nghiệm, và 970 còn trống trong dải này \cite{evomaster611}.

Phương pháp tách làm hai thành phần: bộ dựng kịch bản kiểm thử tạo các bộ ba thao tác (source, transition, follow-up), còn bộ phán quyết đánh giá các bộ ba đã thực hiện thành công. Cách tách này phân biệt được lượt không dựng được bộ ba hợp lệ với lượt có bộ ba nhưng không có vi phạm. Hình~\ref{fig:triple_sequence} trình bày trình tự thực hiện bộ ba và điều kiện kích hoạt Oracle 970.

\begin{figure}[htbp]
  \centering
  % Hình 3.1 (flowchart, thay sequence diagram fig_3_1_triple_sequence): trình tự bộ ba thao tác và kích hoạt Oracle 970.
% Quy tắc nhánh theo Mục 3.3.3 và docs/metamorphic-authz-relations.md: 2xx = cho phép, 403 = từ chối, 401 = hủy.
\begin{tikzpicture}[
  >=Stealth,
  font=\small\linespread{1}\selectfont,
  node distance=6mm and 12mm,
  procbox/.style={rectangle, draw, rounded corners=2pt, fill=gray!10, align=center,
                  text width=5.0cm, minimum height=0.9cm},
  decision/.style={diamond, draw, fill=gray!10, align=center, aspect=2.4, inner sep=1pt, font=\footnotesize},
  voidbox/.style={rectangle, draw, dashed, rounded corners=6pt, align=center, text width=1.9cm,
                  minimum height=0.7cm, font=\footnotesize},
  passbox/.style={rectangle, draw, rounded corners=6pt, align=center, text width=3.2cm, minimum height=0.8cm},
  verdictbox/.style={rectangle, draw, very thick, rounded corners=6pt, align=center, text width=4.4cm, minimum height=0.8cm},
  lbl/.style={font=\footnotesize}
]
  % Cột trái: bộ ba thao tác
  \node[procbox]  (src) {\textbf{Source:} nạn nhân gọi thao tác ghi bằng token $c_v$};
  \node[decision, below=of src] (d1) {2xx?};
  \node[procbox,  below=of d1]  (tr) {\textbf{Transition:} quản trị viên hạ quyền nạn nhân bằng token $c_o$};
  \node[decision, below=of tr]  (d2) {2xx?};
  \node[procbox,  below=of d2]  (fu) {\textbf{Follow-up:} nạn nhân gọi lại source bằng token cũ $c_v$};
  \node[decision, below=of fu]  (d3) {Mã phản hồi?};
  \node[passbox,  below=8mm of d3] (ok) {Phù hợp MR-A1\\(không vi phạm)};

  % Nhánh hủy của cột trái
  \node[voidbox, right=of d1] (v1) {Hủy bộ ba};
  \node[voidbox, right=of d2] (v2) {Hủy bộ ba};
  \node[voidbox, left=of d3]  (v3) {Hủy bộ ba};

  % Cột phải: bước đăng nhập lại và phán quyết
  \node[procbox, text width=4.4cm, right=10mm of d3] (rl) {\textbf{Đăng nhập lại:} chạy lại chuỗi lời gọi bằng token mới $c_v'$ (cần nhận 403) và không kèm xác thực (cần nhận mã khác 403)};
  \node[decision, below=of rl]  (d4)  {Đạt cả hai?};
  \node[verdictbox, below=of d4] (vio) {\textbf{Oracle 970} báo vi phạm MR-A1};
  \node[voidbox, right=9mm of d4] (v4) {Hủy bộ ba};

  \draw[->] (src) -- (d1);
  \draw[->] (d1) -- node[lbl, right] {Có} (tr);
  \draw[->] (d1) -- node[lbl, above] {Không} (v1);
  \draw[->] (tr) -- (d2);
  \draw[->] (d2) -- node[lbl, right] {Có} (fu);
  \draw[->] (d2) -- node[lbl, above] {Không} (v2);
  \draw[->] (fu) -- (d3);
  \draw[->] (d3) -- node[lbl, right] {403} (ok);
  \draw[->] (d3) -- node[lbl, above] {401} (v3);
  \draw[->] (d3) -- node[lbl, above] {2xx} (rl);
  \draw[->] (rl) -- (d4);
  \draw[->] (d4) -- node[lbl, right] {Có} (vio);
  \draw[->] (d4) -- node[lbl, above] {Không} (v4);
\end{tikzpicture}

  \caption{Trình tự thực hiện bộ ba thao tác và kích hoạt Oracle 970}
  \label{fig:triple_sequence}
\end{figure}

\section{Quan hệ biến hình MR-A1}
\label{sec:method_mra1}

MR-A1 là một quan hệ biến hình (Mục~\ref{sec:bg_metamorphic}) giữa hai lần thực hiện một thao tác, trước và sau khi thu hồi quyền. Trong tên MR-A1, chữ A chỉ nhóm quan hệ phân quyền (authorization) và số 1 chỉ quan hệ đầu tiên của nhóm. Các quan hệ trong nhóm có cùng dạng và khác nhau ở cách chủ thể mất quyền: A1 xét trường hợp vai trò bị hạ, còn chuyển quyền sở hữu tài nguyên và quyền hết hạn theo thời gian thuộc các quan hệ khác (Mục~\ref{sec:conclusion_future_work}).

MR-A1 xét một chủ thể $p$, tức tài khoản được kiểm tra, có quyền thực hiện thao tác trước khi bị hạ quyền. Chủ thể này gửi yêu cầu bằng thông tin xác thực $c$ định danh $p$, có thể là token hoặc cookie phiên. Thao tác $a$ là một lời gọi REST API dùng phương thức ghi và có kiểm soát phân quyền, chẳng hạn thao tác cập nhật báo cáo mà chỉ quản trị viên được làm. Trạng thái $s$ là cấu hình phân quyền của hệ thống tại lúc quan sát. Một tài khoản quản trị đổi trạng thái này bằng thao tác $T$ (transition) để hạ vai trò hoặc thu hồi quyền của $p$.

Hàm quan sát $\mathrm{Obs}(a, c, s)$ lấy kết quả từ mã trạng thái HTTP khi chủ thể dùng $c$ gọi thao tác $a$ tại trạng thái $s$ \cite{rfc9110}. Bảng~\ref{tab:obs_status} nêu ba giá trị quan sát và mã trạng thái tương ứng. Giá trị \textsc{allowed} không xác nhận trạng thái dữ liệu nghiệp vụ sau thao tác. Oracle phán quyết dựa trên mã 403. Mã 404 không được tính là lời từ chối, vì tài nguyên không tồn tại cũng trả về mã này. Trong 98 lượt chạy ở Chương~\ref{chap:eval}, không bước follow-up nào nhận mã 404.

Mã \texttt{401} ở bước follow-up nằm ngoài phạm vi phán quyết, vì mã này không phân biệt được hai trường hợp: máy chủ từ chối vì quyền, hay máy chủ không chấp nhận thông tin xác thực. Trên các ứng dụng mục tiêu, nhóm đường dẫn \texttt{/api/*} của pyLoad trả 401 cho phiên không có quyền quản trị, và VAmPI trả 401 khi token không hợp lệ. Bản vá của artifact-keeper thu hồi quyền bằng cách vô hiệu hóa phiên và trả 401 cho token đã bị thu hồi (Mục~\ref{sec:eval_cases}). Nếu coi 401 là \textsc{denied}, một token hết hạn giữa bộ ba sẽ được tính là một lần thu hồi quyền đúng.

\begin{table}[htbp]
  \centering
  \caption{Ánh xạ giá trị quan sát sang mã trạng thái HTTP}
  \label{tab:obs_status}
  \begin{tabularx}{\textwidth}{llX}
    \toprule
    Giá trị & Mã HTTP & Ý nghĩa \\
    \midrule
    \textsc{allowed} & 2xx & Máy chủ chấp nhận lời gọi ở tầng HTTP. \\
    \textsc{denied} & 403 & Máy chủ từ chối lời gọi trên cơ sở phân quyền. \\
    \textsc{inconclusive} & 401 & Phiên hết hạn hoặc chưa xác định được quyền. \\
    \bottomrule
  \end{tabularx}
\end{table}

Gọi $s' = T(s)$ là trạng thái sau khi thực hiện $T$. Khi $c$ được giữ cố định qua hai lần quan sát, MR-A1 có dạng Công thức~\eqref{eq:mra1}.
\begin{equation}
  \mathrm{Obs}(a, c, s) = \textsc{allowed}\ \wedge\ \mathit{Pre}(a, T)\ \implies\ \mathrm{Obs}(a, c, s') = \textsc{denied}.
  \label{eq:mra1}
\end{equation}
Trong đó, $\mathit{Pre}(a, T)$ là điều kiện tiền đề: thao tác $a$ cần đúng quyền mà $T$ thu hồi của chủ thể đang xét. Khi vế trái đúng mà kết quả ở trạng thái $s'$ vẫn là \textsc{allowed}, MR-A1 bị vi phạm và oracle báo lỗi.

Trong thuật ngữ kiểm thử biến hình, lần thực thi nguồn là yêu cầu $a$ gửi bằng $c$ ở trạng thái $s$. Phép biến đổi giữ nguyên yêu cầu và thông tin xác thực, chỉ đổi trạng thái của hệ thống bằng $T$. Lần thực thi tiếp theo gửi lại đúng yêu cầu đó bằng cùng $c$ ở trạng thái $s'$ và không đăng nhập lại. Lỗi cần tái hiện nằm ở phiên cũ: phiên mở trước lúc hạ quyền vẫn giữ quyền đã bị thu hồi.

Điều kiện tiền đề loại các thao tác không phụ thuộc vào quyền bị thu hồi để giảm báo động giả. Chẳng hạn, thao tác cập nhật hồ sơ cá nhân dùng chung cho nhiều vai trò; sau khi tài khoản bị hạ quyền, máy chủ vẫn chấp nhận thao tác này là đúng. Do chuẩn RBAC để ngỏ cách cập nhật quyền trên phiên đang hoạt động \cite{incits359draft}, yêu cầu từ chối ngay sau khi hạ quyền là chính sách riêng của từng hệ thống mục tiêu.

Bảng~\ref{tab:temporal_vs_diff} đối chiếu MR-A1 với nhóm oracle so sánh chéo tài khoản (differential oracle) của Şahin và cộng sự \cite{sahin2024bola}, dùng để phát hiện lỗi BOLA và BFLA. Hình~\ref{fig:mr_a1_state_space} biểu diễn hai trạng thái phân quyền cùng nhánh an toàn và nhánh vi phạm.

\begin{table}[htbp]
  \centering
  \caption{Đối chiếu oracle so sánh chéo tài khoản và quan hệ MR-A1}
  \label{tab:temporal_vs_diff}
  \begin{tabularx}{\textwidth}{@{}l >{\raggedright\arraybackslash}X >{\raggedright\arraybackslash}X@{}}
    \toprule
    Tiêu chí & Oracle so sánh chéo tài khoản & MR-A1 (theo thời gian) \\
    \midrule
    So sánh theo & nhiều tài khoản tại một thời điểm & một tài khoản trước và sau khi hạ quyền \\
    Thông tin xác thực & khác nhau theo tài khoản & giữ cố định và dùng lại \\
    Nội dung kiểm tra & tài khoản không có quyền vẫn truy cập được & quyền đã thu hồi vẫn được chấp nhận \\
    \bottomrule
  \end{tabularx}
\end{table}

\begin{figure}[htbp]
  \centering
  \begin{tikzpicture}[
  >=Stealth,
  font=\small\linespread{1}\selectfont,
  state/.style={rectangle, draw, rounded corners=4pt, fill=gray!10, align=center,
                text width=5.4cm, minimum height=1.3cm},
  safe/.style={rectangle, draw, rounded corners=6pt, align=center, text width=4.6cm, minimum height=1.2cm},
  vuln/.style={safe, very thick},
  lbl/.style={font=\footnotesize, align=left}
]
  \node[state] (s0) at (0, 0)    {\textbf{Trạng thái $s$}\\vai trò quản trị, token $c$};
  \node[state] (s1) at (0, -3.0) {\textbf{Trạng thái $s' = T(s)$}\\vai trò thấp hơn, vẫn dùng token $c$};
  \node[safe] (ok)  at (-3.3, -6.0) {\textbf{Phù hợp MR-A1}\\follow-up nhận 403};
  \node[vuln] (bad) at (3.3, -6.0)  {\textbf{Vi phạm MR-A1}\\follow-up nhận 2xx,\\Oracle 970 báo lỗi};

  \draw[->] ([yshift=3mm]s0.east) -- ++(9mm, 0) |- node[lbl, pos=0.25, right] {source $a$\\nhận 2xx} ([yshift=-3mm]s0.east);
  \draw[->] (s0) -- node[lbl, right] {transition $T$ hạ vai trò,\\token $c$ giữ nguyên} (s1);
  \draw[->] (s1) -- node[lbl, above left] {403} (ok);
  \draw[->] (s1) -- node[lbl, above right] {2xx} (bad);
\end{tikzpicture}

  \caption{Hai trạng thái phân quyền trong MR-A1}
  \label{fig:mr_a1_state_space}
\end{figure}

\section{Bộ ba thao tác và điều kiện tiền đề}
\label{sec:method_triple}

Bộ dựng kịch bản kiểm thử chuyển quan hệ MR-A1 thành ca kiểm thử chạy được trên REST API, gọi là kịch bản kiểm thử MR-A1. Đầu vào của bộ dựng là các ca kiểm thử trong kho lưu trữ, gồm cả những ca nhận mã 403 do bước sinh phản hồi 403 có sẵn của EvoMaster tạo ra \cite{evomaster611}.

\subsection{Cấu trúc bộ ba thao tác}
\label{subsec:method_triple_structure}

Kịch bản kiểm thử MR-A1 gồm các lời gọi khởi tạo và bộ ba thao tác chính, chạy theo thứ tự sau:
\begin{enumerate}
    \item Khởi tạo (setup): các lời gọi trong ca kiểm thử ban đầu của nạn nhân, chạy trước thao tác gốc để tạo tài nguyên cần thao tác. Bước này không bắt buộc nếu tài nguyên đã có sẵn trong hệ thống.
    \item Source (thao tác gốc): tài khoản nạn nhân, trước khi bị hạ quyền, gọi một thao tác ghi có kiểm soát phân quyền và nhận mã \texttt{2xx}.
    \item Transition (thao tác chuyển trạng thái): tài khoản quản trị hạ vai trò hoặc thu hồi quyền tương ứng của nạn nhân.
    \item Follow-up (thao tác kiểm chứng): nạn nhân gọi lại thao tác source với cùng tham số và cùng token cũ.
\end{enumerate}

\begin{table}[htbp]
  \centering
  \caption{Kịch bản kiểm thử MR-A1 trên reports-api}
  \label{tab:triple_example}
  \begin{tabularx}{\textwidth}{@{}c l >{\raggedright\arraybackslash}X c c@{}}
    \toprule
    \# & Tài khoản & Lời gọi & Bản lỗi & Bản vá \\
    \midrule
    1 & \texttt{alice} & Khởi tạo: \texttt{POST /reports} & 201 & 201 \\
    2 & \texttt{alice} & Source: \texttt{PUT /reports/\{id\}} trên báo cáo vừa tạo & 200 & 200 \\
    3 & \texttt{superadmin} & Transition: \texttt{PUT /users/alice/role}, đặt \texttt{role} thành \texttt{VIEWER} & 200 & 200 \\
    4 & \texttt{alice} & Follow-up: lặp lại lời gọi 2 với cùng token và cùng tham số & 200 & 403 \\
    \bottomrule
  \end{tabularx}
\end{table}

Follow-up dùng lại được token cũ vì EvoMaster đăng nhập mỗi tài khoản một lần ở đầu ca kiểm thử, rồi dùng token hoặc cookie phiên đó cho mọi thao tác sau.

Bảng~\ref{tab:triple_example} là một kịch bản kiểm thử MR-A1 do công cụ dựng trên reports-api (Mục~\ref{subsec:impl_reports}). Tài khoản \texttt{alice} có vai trò quản trị, còn \texttt{superadmin} là tài khoản quản trị hạ quyền \texttt{alice} xuống \texttt{VIEWER}.

\subsection{Điều kiện tiền đề và cách chọn source}
\label{subsec:method_guard}

Bộ dựng kiểm tra điều kiện tiền đề ngay khi chọn source. Source phải là thao tác ghi, xác định theo phương thức HTTP (\texttt{PUT}, \texttt{PATCH}, \texttt{POST}), và khác đường dẫn với thao tác chuyển trạng thái. Bộ dựng chọn một thao tác ghi làm source khi kho lưu trữ có đồng thời hai quan sát sau:
\begin{enumerate}
    \item Thao tác này từng nhận mã 403 ít nhất một lần. Ngoài các lần 403 xuất hiện trong pha tìm kiếm, bước sinh phản hồi 403 tạo thêm quan sát này bằng cách gửi lại một ca kiểm thử thành công bằng tài khoản khác.
    \item Tài khoản nạn nhân từng thực hiện thành công thao tác này với mã \texttt{2xx}.
\end{enumerate}

Cặp phản hồi 403 và \texttt{2xx} này chỉ là dấu hiệu ban đầu, vì lời từ chối có thể đến từ quyền sở hữu tài nguyên hoặc một điều kiện khác. Phân loại theo phương thức cho phép source là một lời gọi \texttt{POST} chỉ đọc dữ liệu. Trên pyLoad, source \texttt{POST /json/status} trả trạng thái máy chủ và bị chặn bằng quyền \texttt{LIST}, quyền mà tài khoản mất khi bị hạ.

Bộ dựng ưu tiên \texttt{PUT} và \texttt{PATCH}: follow-up lặp lại source trên cùng tài nguyên, và hai phương thức này thường gọi lại được nhiều lần. Bộ dựng chỉ xét \texttt{POST} khi không thao tác \texttt{PUT} hay \texttt{PATCH} nào có lời gọi thành công bằng tài khoản đã đăng nhập. \texttt{DELETE} không được chọn làm source, vì sau khi xóa, follow-up không còn tài nguyên để gọi lại.

\subsection{Bộ ba bị hủy và bước đăng nhập lại}
\label{subsec:method_void}

Bộ ba bị hủy và loại khỏi tập đánh giá khi source hoặc transition trả mã ngoài nhóm \texttt{2xx}, hoặc khi follow-up không ghi được mã trạng thái hay nhận mã \texttt{401}. Sau mỗi lần hủy, bộ dựng chọn tổ hợp mới theo các luật ở Mục~\ref{subsec:method_builder_algo}. Với các bộ ba còn lại, mã của follow-up quyết định bước tiếp theo: 403 nghĩa là bộ ba phù hợp MR-A1, còn \texttt{2xx} dẫn tới bước đăng nhập lại. Với mã khác, như 404 hay 500, bộ ba được giữ và không bị báo vi phạm; các lượt chạy của khóa luận không gặp trường hợp này.

Trong một số lần thử, mọi tài khoản quản trị đều bị từ chối khi hạ quyền, dù source của nạn nhân vẫn thành công. Tình huống này xảy ra khi pha tìm kiếm đã hạ quyền các tài khoản quản trị nhưng chưa hạ quyền tài khoản nạn nhân. Bộ dựng khi đó dùng quyền quản trị còn hiệu lực của tài khoản nạn nhân để khôi phục vai trò cho tài khoản quản trị rồi thử lại. Lần khôi phục là một ca kiểm thử riêng, chạy trước bộ ba mới: nạn nhân gọi chính thao tác chuyển trạng thái để nâng quyền cho tài khoản quản trị. Ca kiểm thử này không tự tạo ra cảnh báo, vì bộ ba dựng sau đó vẫn phải qua mọi điều kiện hủy và bước đăng nhập lại. Trên pyLoad, ở mọi lượt chạy, các tài khoản quản trị đều bị từ chối khi hạ quyền nạn nhân \texttt{pyload}; \texttt{pyload} nâng quyền \texttt{alice} trở lại, rồi \texttt{alice} hạ quyền \texttt{pyload} để dựng bộ ba. Bước khôi phục chỉ đảo được việc hạ quyền: tài khoản đã bị xóa trong pha tìm kiếm không phục hồi được.

Khi follow-up nhận \texttt{2xx}, bộ dựng chạy thêm bước đăng nhập lại trước khi chuyển bộ ba sang bộ phán quyết. Bước này kiểm tra transition có thu hồi quyền hay không. Tên endpoint chưa đủ để biết: một endpoint có tên gợi ý hạ quyền vẫn có thể nâng quyền, tùy giá trị vai trò gửi kèm.

Ở bước đăng nhập lại, bộ dựng dùng token mới của tài khoản nạn nhân để chạy lại chuỗi lời gọi của nạn nhân trong ca kiểm thử gốc. Chuỗi này gồm các lời gọi khởi tạo dữ liệu và thao tác source. Do EvoMaster dừng chuỗi tại lời gọi tạo tài nguyên đầu tiên thất bại, bộ dựng đọc mã trạng thái của lời gọi cuối cùng được thực thi. Mã 403 tại lời gọi này cho thấy token mới đã mất quyền ghi.

Khi nhận mã 403, bộ dựng chạy cùng chuỗi một lần nữa, không kèm thông tin xác thực. Nếu lần gọi này cũng nhận 403, mã 403 nhận được bằng token mới chưa chắc do quyền bị thu hồi. Bộ ba chỉ được giữ khi lần gọi không kèm xác thực trả về một mã khác 403, chẳng hạn 401. Mọi trường hợp còn lại, kể cả khi token mới vẫn nhận \texttt{2xx}, đều làm bộ ba bị hủy.

Với kịch bản ở Bảng~\ref{tab:triple_example} trên biến thể có lỗ hổng, bước đăng nhập lại chạy lại lời gọi 1 và 2 bằng token mới của \texttt{alice}. Lời gọi tạo báo cáo nhận 403 nên chuỗi dừng tại đó; lượt gửi không kèm thông tin xác thực nhận 401. Bộ ba được giữ và Oracle 970 báo vi phạm.

Bước đăng nhập lại chạy sau follow-up. Ở hệ thống giữ một phiên cho mỗi tài khoản, lần đăng nhập mới sẽ vô hiệu hóa token cũ mà follow-up cần. Token mới chỉ có được qua lần đăng nhập ở đầu ca kiểm thử, nên bước này chạy trong một ca kiểm thử riêng.

Công cụ hộp đen không quan sát trực tiếp được điều kiện tiền đề $\mathit{Pre}(a, T)$, vì quy tắc phân quyền nằm trong logic nội bộ của ứng dụng. Bộ dựng ước lượng điều kiện này bằng dấu hiệu HTTP ở hai thời điểm: khi chọn source và ở bước đăng nhập lại. Ở thời điểm thứ hai, mã 403 có thể đến từ một lời gọi khởi tạo thay vì từ chính source. Một cảnh báo của Oracle 970 là vi phạm MR-A1 khi hai giả định cùng đúng: mã 403 phản ánh quyết định phân quyền, và phiên mới mang quyền hiện hành. Nếu giả định thứ hai sai, tức token mới vẫn mang quyền cũ, bước đăng nhập lại có thể loại nhầm một bộ ba chứa lỗ hổng.

\subsection{Thuật toán dựng kịch bản kiểm thử MR-A1}
\label{subsec:method_builder_algo}

Thuật toán~\ref{alg:build_stateful_triple} duyệt các thao tác chuyển trạng thái theo thứ tự xếp hạng và ghép mỗi thao tác với các thao tác ghi ứng viên trong kho lưu trữ. Mỗi tổ hợp chạy trong giới hạn số lần thử, kèm bước khôi phục quyền khi thao tác chuyển trạng thái bị từ chối.

Trong thuật toán, $c_v$ và $c_o$ là thông tin xác thực của tài khoản nạn nhân $v$ và tài khoản quản trị. Dòng 4 ước lượng $\mathit{Pre}(a, T)$ bằng hai quan sát khi chọn source. Ở dòng 8, bộ dựng áp các điều kiện hủy cho source, transition và follow-up, gồm cả bước đăng nhập lại khi follow-up nhận \texttt{2xx}. Bước khôi phục quyền nằm ở dòng 11; lần thử kế tiếp cho biết bước này có thành công hay không.

\begin{algorithm}[htbp]
\caption{Dựng kịch bản kiểm thử MR-A1}
\label{alg:build_stateful_triple}
\begin{algorithmic}[1]
\Require kho lưu trữ $\mathcal{A}$; giới hạn lần thử $K$, $N$; giới hạn khôi phục $R$
\Ensure tập $\mathcal{I}$ các bộ ba đã thực hiện
\State $n \gets 0$; $\mathcal{I} \gets \emptyset$
\For{mỗi thao tác chuyển trạng thái $T$ theo thứ tự xếp hạng}
    \State $k \gets 0$; $r \gets 0$
    \For{mỗi thao tác ghi $a \neq T$ thỏa hai quan sát ở Mục~\ref{subsec:method_guard}}
        \While{$k < K$ \textbf{và} $n < N$}
            \State $I \gets [\text{khởi tạo của } a,\ a(c_v),\ T(c_o, v),\ a(c_v)]$
            \State chạy $I$
            \If{$I$ không bị hủy}
                \State $\mathcal{I} \gets \mathcal{I} \cup \{I\}$; \textbf{dừng}
            \ElsIf{$T$ nhận mã 401 hoặc 403 \textbf{và} $r < R$}
                \State dùng $c_v$ khôi phục quyền cho tài khoản quản trị; $r \gets r + 1$
            \EndIf
            \State $k \gets k + 1$; $n \gets n + 1$
        \EndWhile
    \EndFor
\EndFor
\State \Return $\mathcal{I}$
\end{algorithmic}
\end{algorithm}

Ba trần giới hạn chi phí chạy: tổng $N = 20$ lần thử trong cả pha, $K = 5$ lần thử và $R = 3$ lần khôi phục quyền cho mỗi thao tác chuyển trạng thái. Thuật toán~\ref{alg:build_stateful_triple} là bản rút gọn: trong cài đặt, trần $K$ không áp dụng cho thao tác chuyển trạng thái đã dựng được bộ ba và cho ứng viên cuối của danh sách. Với mỗi thao tác ghi ứng viên, bộ ba hợp lệ đầu tiên được đưa vào tập kết quả cho bộ phán quyết, kể cả khi follow-up nhận 403. Bộ dựng sau đó chuyển sang ứng viên kế tiếp; một lượt chạy có thể cho nhiều bộ ba.

Khi một lần thử thất bại, bộ dựng chọn bước kế tiếp dựa trên mã trạng thái của thao tác chuyển trạng thái (Bảng~\ref{tab:retry_rules}). Tài khoản quản trị nào từng gọi thành công thao tác chuyển trạng thái trong kho lưu trữ được thử trước. Các luật bỏ qua giúp bộ dựng không tiêu lượt thử vào những lần đã biết sẽ thất bại.

\begin{table}[htbp]
  \centering
  \caption{Luật thử lại và bỏ qua của bộ dựng}
  \label{tab:retry_rules}
  \begin{tabularx}{\textwidth}{@{}>{\raggedright\arraybackslash}X >{\raggedright\arraybackslash}X@{}}
    \toprule
    Tình huống & Bước tiếp theo \\
    \midrule
    Transition nhận mã \texttt{4xx} hoặc \texttt{5xx} khác 401, 403, 404 & Thử giá trị hạ quyền kế tiếp với cùng tài khoản quản trị \\
    Transition nhận 401 hoặc 403 & Thử tài khoản quản trị kế tiếp, rồi khôi phục quyền khi đã thử hết \\
    Transition nhận \texttt{2xx} & Không dùng lại nạn nhân này cho thao tác chuyển trạng thái đang thử \\
    Transition từ chối mọi lần thử trên một nạn nhân & Không thử lại cặp thao tác chuyển trạng thái và nạn nhân này trong pha \\
    Transition nhận 404 sau source thành công, khi chưa dựng được bộ ba nào & Bỏ thao tác chuyển trạng thái này, chuyển sang ứng viên kế \\
    Đường dẫn của ứng viên và đường dẫn cha chỉ nhận 404 trong kho lưu trữ & Loại ứng viên trước khi thử \\
    \bottomrule
  \end{tabularx}
\end{table}

\section{Tự động nhận diện và xếp hạng thao tác chuyển trạng thái}
\label{sec:method_ranking}

Số lần thử của Thuật toán~\ref{alg:build_stateful_triple} có giới hạn, nên bộ dựng thử các thao tác chuyển trạng thái theo thứ tự xếp hạng. Mô hình Naive Bayes lập thứ tự này bằng cách chấm điểm từng thao tác ghi trong đặc tả OpenAPI. Hình~\ref{fig:transition_inference} tóm tắt quy trình nhận diện.

\begin{figure}[htbp]
  \centering
  \begin{tikzpicture}[
  >=Stealth,
  font=\small\linespread{1}\selectfont,
  node distance=8mm and 7mm,
  procbox/.style={rectangle, draw, rounded corners=2pt, fill=gray!10, align=center,
                  text width=4.2cm, minimum height=1.5cm},
  outbox/.style={rectangle, draw, rounded corners=2pt, align=left,
                 text width=7.4cm, minimum height=1.5cm}
]
  \node[procbox] (spec) {\textbf{Đặc tả OpenAPI}\\các thao tác ghi};
  \node[procbox, right=of spec] (tok) {\textbf{Tách đặc trưng}\\đường dẫn, tên trường\\và phương thức HTTP};
  \node[procbox, right=of tok] (nb) {\textbf{Naive Bayes}\\$\mathrm{score}(op)$ bằng tổng\\trọng số logarit};
  \node[procbox, below=of nb] (rank) {\textbf{Xếp hạng}\\giữ thao tác có\\$\mathrm{score}(op) > 0$};
  \node[outbox, left=of rank] (out) {\textbf{Đầu vào cho bộ dựng}\\1. Thao tác hạ quyền theo thứ tự điểm\\2. Trường vai trò và giá trị hạ quyền\\3. Trường định danh tài khoản nạn nhân};

  \draw[->] (spec) -- (tok);
  \draw[->] (tok) -- (nb);
  \draw[->] (nb) -- (rank);
  \draw[->] (rank) -- (out);
\end{tikzpicture}

  \caption{Quy trình nhận diện thao tác hạ quyền từ đặc tả OpenAPI}
  \label{fig:transition_inference}
\end{figure}

Naive Bayes là bộ phân loại xác suất dựa trên định lý Bayes, giả định các đặc trưng độc lập có điều kiện khi biết nhãn. Mô hình học từ 10.381 thao tác ghi của 230 REST API. Một thao tác mang nhãn chuyển trạng thái khi nhắm vào một tài khoản đã tồn tại hoặc chính người gọi, đổi vai trò, quyền hay cờ quản trị của tài khoản đó, và không thuộc loại chỉ cấp thêm quyền. Các thao tác còn lại, kể cả thao tác tạo tài khoản mới hay khóa tài khoản, mang nhãn thông thường.

Hai mô hình ngôn ngữ lớn gán nhãn khởi đầu, mỗi mô hình một lượt độc lập. Sau đó 602 thao tác được soát tay: 66 thao tác hai lượt gán lệch nhau, 165 trong 173 thao tác hai lượt cùng gán nhãn chuyển trạng thái và khớp nhau hoàn toàn, và 371 thao tác hai lượt cùng gán nhãn thông thường, chọn ngẫu nhiên. Mẫu ngẫu nhiên này không có thao tác chuyển trạng thái nào. Thao tác chưa được soát giữ nhãn mà hai lượt gán cùng cho. Tập nhãn cuối có 209 thao tác chuyển trạng thái.

Mô hình dùng trong công cụ được học từ toàn bộ tập này và được chấm trên một tập đặc tả riêng (Mục~\ref{sec:eval_ranking}). Ngoài bảng trọng số dùng để xếp hạng thao tác, mô hình còn các bảng trọng số riêng để chọn trường vai trò và trường định danh tài khoản nạn nhân. Giá trị hạ quyền phụ thuộc kiểu của trường vai trò (Bảng~\ref{tab:role_values}); chỉ trường kiểu liệt kê và kiểu chuỗi dùng mô hình. Bộ dựng loại những ứng viên không đề xuất được giá trị hạ quyền nào, chẳng hạn vì không tìm thấy trường vai trò. Trên pyLoad, lần hạ quyền có hiệu lực nhờ giá trị \texttt{role} do pha tìm kiếm để lại; công cụ không tự suy ra giá trị này (Mục~\ref{sec:eval_cases}).

\begin{table}[htbp]
  \centering
  \caption{Cách chọn giá trị hạ quyền theo kiểu của trường vai trò}
  \label{tab:role_values}
  \begin{tabularx}{\textwidth}{@{}l >{\raggedright\arraybackslash}X l@{}}
    \toprule
    Kiểu trường & Giá trị hạ quyền & Ví dụ \\
    \midrule
    Liệt kê & Các giá trị khai báo, xếp theo mức mô hình coi là quyền thấp, trừ vai trò hiện tại của nạn nhân & reports-api: \texttt{VIEWER} \\
    Logic & Giá trị đảo & artifact-keeper: \texttt{is\_admin = false} \\
    Số & Lần lượt \texttt{1}, \texttt{0}, \texttt{2}, thứ tự cố định trong mã & pyLoad: \texttt{permission = 1} \\
    Chuỗi & Danh sách giá trị quyền thấp của mô hình & \\
    Mảng & Mảng rỗng & \\
    \bottomrule
  \end{tabularx}
\end{table}

Với mỗi đặc trưng $t$, mô hình tính một trọng số logarit theo tần suất của $t$ ở hai lớp. Điểm của một thao tác $op$ được tính theo Công thức~\eqref{eq:nb_weight}:
\begin{equation}
  w(t) = \ln\frac{d_T + 1}{N_T + 2} - \ln\frac{d_N + 1}{N_N + 2}, \qquad
  \mathrm{score}(op) = \ln\frac{N_T}{N_N} + \sum_{t \in op} w(t).
  \label{eq:nb_weight}
\end{equation}
Trong đó, $d_T$ và $d_N$ là số thao tác thuộc lớp chuyển trạng thái và lớp thông thường có chứa $t$, còn $N_T$ và $N_N$ là tổng số thao tác của hai lớp. Thao tác có $\mathrm{score}(op) > 0$ được coi là ứng viên chuyển trạng thái. Trường hợp hai thao tác có điểm số bằng nhau, bộ dựng ưu tiên theo thứ tự xuất hiện ban đầu của chúng trong đặc tả OpenAPI. Điểm số chỉ dùng để xếp hạng các ứng viên và chưa được hiệu chuẩn để đọc như một xác suất hậu nghiệm.

Đặc trưng của một thao tác gồm ba nhóm: các từ trong đường dẫn (\texttt{p:}), các từ trong tên tham số truy vấn và tên trường của thân yêu cầu (\texttt{f:}), và phương thức HTTP (\texttt{v:}). Tên được tách thành từ tại các ký tự không phải chữ cái hay chữ số, và tại chỗ đổi chữ hoa trong tên kiểu camelCase, rồi chuyển về chữ thường; từ toàn chữ số và từ ngắn hơn 4 ký tự bị bỏ. Mô hình chỉ giữ những đặc trưng xuất hiện trong thao tác chuyển trạng thái của ít nhất hai nhà cung cấp API, để mô hình không khớp theo cách đặt tên của riêng một công ty. Đặc trưng được giữ vẫn có thể mang trọng số âm, chẳng hạn \texttt{v:POST}. Tên tham số đường dẫn đã có trong \texttt{p:}, nên không được lấy lại vào \texttt{f:}; nếu lấy lại, cùng một từ được cộng hai lần.

Bảng~\ref{tab:nb_example} tính điểm của thao tác \texttt{PUT /users/\{username\}/role} trên reports-api. Tiên nghiệm âm vì chỉ 209 trong 10.381 thao tác của tập học mang nhãn chuyển trạng thái. Năm đặc trưng có mặt đưa điểm lên 5,63. Đây là thao tác duy nhất của reports-api có điểm dương. Thao tác \texttt{PUT /reports/\{id\}} chỉ có đặc trưng \texttt{v:PUT} trong mô hình: các từ \texttt{reports}, \texttt{title}, \texttt{classification} không có trong mô hình, còn \texttt{id} ngắn hơn 4 ký tự. Điểm của thao tác này là $-3{,}885 + 0{,}250 = -3{,}64$.

\begin{table}[htbp]
  \centering
  \caption{Điểm của thao tác \texttt{PUT /users/\{username\}/role} trên reports-api}
  \label{tab:nb_example}
  \begin{tabular}{@{}l r r@{}}
    \toprule
    Thành phần & Trọng số & Cộng dồn \\
    \midrule
    Tiên nghiệm & $-3{,}885$ & $-3{,}885$ \\
    \texttt{p:users} & $+2{,}135$ & $-1{,}750$ \\
    \texttt{p:username} & $+2{,}449$ & $+0{,}698$ \\
    \texttt{p:role} & $+2{,}755$ & $+3{,}453$ \\
    \texttt{f:role} & $+1{,}925$ & $+5{,}378$ \\
    \texttt{v:PUT} & $+0{,}250$ & $+5{,}628$ \\
    \bottomrule
  \end{tabular}
\end{table}

Bộ dựng và bộ phán quyết dùng chung quy tắc $\mathrm{score}(op) > 0$. Bộ phán quyết chỉ nhận lời gọi ở giữa làm thao tác chuyển trạng thái khi lời gọi đó thỏa quy tắc, và bỏ qua source thỏa quy tắc. Vì vậy bộ dựng không dựng bộ ba nào mà bộ phán quyết từ chối. Ngưỡng tính cả tiên nghiệm. Nếu bỏ tiên nghiệm, \texttt{PUT /reports/\{id\}} đạt $+0{,}250$ nhờ \texttt{v:PUT}, bị coi là thao tác chuyển trạng thái và bị loại khỏi vai trò source; Oracle 970 khi đó không bao giờ báo vi phạm trên reports-api.

Giả định độc lập của mô hình bị vi phạm ở các từ thường đi cùng nhau, như \texttt{user}, \texttt{username}, \texttt{email}. Trong Bảng~\ref{tab:nb_example}, hai đặc trưng về người dùng góp 4,58 trong 5,63 điểm. Điểm số vì vậy quá tự tin: điểm 5,63 ứng với hậu nghiệm 0,996, trong khi trên tập chấm chỉ 101 trong 624 thao tác có điểm dương mang nhãn chuyển trạng thái (Mục~\ref{sec:eval_ranking}).

\endgroup
